\subsection{借助可测等价关系的测度分解}

设 T 是一个 Polish 局部紧空间， \(\mu\) 是 T 上的正测度， R 是 T 中的一个 \(\mu\) -可测等价关系。则（No.~4, 命题 2）存在一个 Polish 局部紧空间 B 和 T 到 B 中的一个 \(\mu\) -可测映射 p ，使得关系 \(p(x) = p(y)\) 等价于 \(R\{x, y\}\) 。每个测度 \(\nu\) （在 p 下为 \(\mu\) 的伪像）（No.~2）都称为 \(\mu\) 按关系 R 的商测度；若 \(b \mapsto \lambda_b\) 是 \(\mu\) 关于测度 \(\nu\) 的一个分解，我们就说 \(b \mapsto \lambda_b\) 是 \(\mu\) 借助关系 R 的一个分解。由于 p 与诸 \(\lambda_b\) 的性质，每个测度 \(\lambda_b\) 都集中于一个关于 R 的等价类上，且若 \(b \neq c\) ，测度 \(\lambda_b\) 与 \(\lambda_c\) 集中于不同的类上。

空间 B 、映射 p 以及 B 上的伪像测度 \(\nu\) 一般可以用无穷多种方式选取。尽管如此， \(\mu\) 借助 R 的各个分解都可以由其中之一导出，这是下述定理的一个推论：

定理 4. --- 设 T 是一个 Polish 局部紧空间， \(\mu\) 是 T 上的正测度， R 是 T 中的一个 \(\mu\) -可测等价关系。设 B 、 B' 是两个 Polish 局部紧空间， \(p, p'\) 是 T 到 B 、 B' 中的两个 \(\mu\) -可测映射，使得 R\{x,y\} 等价于 \(p(x) = p(y)\) 且等价于 \(p'(x) = p'(y)\) 。设 \(\nu, \nu'\) 分别是 \(\mu\) 在 \(p, p'\) 下的伪像测度；设 \(b \mapsto \lambda_b\) 、 \(b' \mapsto \lambda_{b'}'\) 分别是 \(\mu\) 关于 \(\nu, \nu'\) 的分解。

在这些条件下， B （相应地， B' ）中存在一个关于 \(\nu\) （相应地， \(\nu'\) ）可忽略的集 N （相应地， N' ），以及 B - N 到 B' - N' 上的一个双射 f ，具有下列性质：

\begin{enumerate}
\def\labelenumi{\alph{enumi})}
\item
  映射 \(f\) （在 B 中几乎处处有定义）是 \(\nu\) -可测的，其逆映射 \(f'\) 是 \(\nu'\) -可测的； \(\nu\) （相应地， \(\nu'\) ）在 \(f\) （相应地， \(f'\) ）下的每个伪像测度都等价于 \(\nu'\) （相应地， \(\nu\) ）。
\item
  对每个 \(b \in B - N\) ， T 上的测度 \(\lambda_{f(b)}^{\prime}\) 具有 \(r(b)\lambda_{b}\) 的形式，其中 \(r(b) \neq 0\) 且 r 局部 \(\nu\) -可积。
\end{enumerate}

为证明 \(a\) ），我们可以局限于 \(\nu\) 与 \(\nu'\) 都是有界测度的情形（Ch. V, §5, No.~6, 命题 11）。设 \(\mathrm{N}_0 = \mathrm{B} - p(\mathrm{T})\) ， \(\mathrm{N}_0' = \mathrm{B}' - p'(\mathrm{T})\) ；我们知道 \(\mathrm{N}_0\) （相应地， \(\mathrm{N}_0'\) ）关于 \(\nu\) （相应地， \(\nu'\) ）是可忽略的（No.~2）。存在一个双射 \(f\) （ \(\mathrm{B} - \mathrm{N}_0\) 到 \(\mathrm{B}' - \mathrm{N}_0'\) 上的双射），它由 \(f(p(t)) = p'(t)\) （对一切 \(t \in T\) ）定义；设 \(f'\) 为 \(f\) 的逆映射，从而 \(f'(p'(t)) = p(t)\) 。对每个子集 \(M\) （ \(B\) 的子集），关系« M 是 \(\nu\) -可测的»等价于« \(\overline{p}^{-1}(M)\) 是 \(\mu\) -可测的»，也就是等价于« \(p'(f(M))\) 是 \(\mu\) -可测的»，从而最终等价于« \(f(M)\) 是 \(\nu'\) -可测的»（Ch. V, §6, No.~2, 命题 3 的推论）。于是我们看到， \(f\) （相应地， \(f'\) ）把每个 \(\nu\) -可测（相应地， \(\nu'\) -可测）集变成 \(\nu'\) -可测（相应地， \(\nu\) -可测）集；由于 B 与 \(B'\) 可度量化且具有可数基，由此得出 \(f\) 与 \(f'\) 可测（Ch. IV, §5, No.~5, 定理 4）。此外，若 \(M \subset B\) 是 \(\nu\) -可忽略的，则 \(\overline{p}^{-1}(M) = p'(f(M))\) 是 \(\mu\) -可忽略的，因此 \(f(M)\) 是 \(\nu'\) -可忽略的（Ch. V, §6, No.~2, 命题 2 的推论 2）；类似地， \(f'\) 把每个 \(\nu'\) -可忽略集变成 \(\nu\) -可忽略集。因此， \(\nu\) 在 \(f\) 下的像（它有定义，因为 \(\nu\) 有界，这蕴含 \(f\) 是 \(\nu\) -逆紧的）等价于 \(\nu'\) ，而 \(\nu'\) 在 \(f'\) 下的像等价于 \(\nu\) （Ch. V, §5, No.~6, 命题 10）。剩下要证明 \(b\) ）。根据 No.~3 的定理 2，我们可以局限于 \(\nu' = f(\nu)\) 的情形。由于 \(\mu = \int \lambda_b' d\nu'(b')\) ，对每个函数 \(h \in \mathcal{K}(T)\) ，有

\[
\int h (t) d \mu (t) = \int d \nu^ {\prime} (b ^ {\prime}) \int h (t) d \lambda_ {b ^ {\prime}} ^ {\prime} (t) = \int d \nu (b) \int h (t) d \lambda_ {f (b)} ^ {\prime} (t)
\]

（Ch. V, §3, No.~4, 定理 1 与 §6, No.~2, 定理 1）；换言之， \(\mu = \int \lambda_{f(b)}' d\nu(b)\) 。但由于也有 \(\mu = \int \lambda_b d\nu(b)\) ，且对每个 \(b \in B - N_0\) ， \(\lambda_b\) 与 \(\lambda_{f(b)}'\) 都由 \(\overline{p}^{-1}(b)\) 承载，故 No.~3 的定理 2 蕴含 \(\lambda_b = \lambda_{f(b)}'\) （对几乎每个 \(b \in B - N_0\) ，从而对几乎每个 \(b \in B\) ）。因此，取 \(N\) 为 \(N_0\) 与点 \(b \in B\) （使 \(\lambda_b \neq \lambda_{f(b)}'\) 者）所成的集的并，定理 4 的条件便得到验证。

\appendix
